\documentclass[11pt]{article}
\usepackage[margin=0.85in]{geometry}
\usepackage[T1]{fontenc}
\usepackage{lmodern}
\pagestyle{empty}
\setlength{\parindent}{0pt}
\setlength{\parskip}{4pt}
\begin{document}

\textbf{DUAN Fang}\\
Huaqiao University, School of Information Science and Engineering, Xiamen, China\\
\texttt{duan.fang@hqu.edu.cn}

\vspace{6pt}
\today

\vspace{10pt}
To the Editors,\\
\emph{Journal of Economic Dynamics and Control}

\vspace{6pt}
Dear Editors,

I submit the enclosed manuscript, ``Anatomy of a Smooth Transition: A Task-Based Macroeconomic Model of AGI Displacement with Fiscal and Social Safety Nets'', for consideration as a research article in the \emph{Journal of Economic Dynamics and Control}.

Artificial general intelligence differs from earlier automation in one respect that matters for dynamics: it takes part in the creation of new tasks itself, so the mechanism that historically reinstated displaced labor can shut down endogenously. The transition question then ceases to be how fast output grows and becomes whether the high-employment equilibrium still exists, and how much horizon policy retains before the branch disappears.

The paper contributes on three margins. First, it identifies the collapse as a boundary equilibrium bifurcation of the piecewise-smooth fast subsystem, derives a closed-form critical reinstatement rate, and makes that rate endogenous to public employment, workfare and retraining, so that policy moves the threshold instead of the outcome. Second, it uses geometric singular perturbation theory to characterize the finite-rate lag between the slow fiscal layer and the fast employment layer, which is where most of the transitional cost is incurred. Third, it places the demand side (nominal rigidity, the zero lower bound, debt-deflation), the fiscal side (a capital tax whose base contracts as the labor share vanishes) and an independent social-instability channel inside one numerically solvable switching system, and reports the policy ranking that follows.

The numerical claims are stated as what they are. The loss-minimizing design on the reported grid is a workfare-dominant corner financed by a minimal capital tax; the minimizer is grid-dependent rather than identified, and the variant that retains a buffer is reported because a frontier-growth shock experiment shows that the buffer, not the last hundredth of loss, decides survival. Sufficiency is conditional on the pace of automation: at a frontier growth of five percent the floor is already breached. Section 5 states which magnitudes are anchored to data and which are declared settings, and the concluding section lists the two objects that still lack an external anchor---the endogenous reinstatement coefficients and the reduced-form social channel.

The paper's central object is dynamics and computation: a switching system with two time scales, solved numerically over a grid and analyzed through bifurcation theory and singular perturbation. Much of the current debate on automation and the labor share is conducted in comparative statics, where this paper is complementary rather than duplicative, and I hope the journal's readership will find the bifurcation account of interest.

The manuscript is original, is not under consideration elsewhere, and has not been published previously. I have no competing interests to declare. A generative AI assistant was used for language editing only; the author reviewed and edited all content and takes full responsibility for it. The simulation code and parameter files that generate every reported number are available from the corresponding author on request.

\vspace{6pt}
Sincerely,\\
DUAN Fang

\end{document}
